%% Согласно ГОСТ Р 7.0.11-2011:
%% 5.3.3 В заключении диссертации излагают итоги выполненного исследования, рекомендации, перспективы дальнейшей разработки темы.
%% 9.2.3 В заключении автореферата диссертации излагают итоги данного исследования, рекомендации и перспективы дальнейшей разработки темы.
\begin{enumerate}
	\item Сделано обобщение понятия моста Эйнштейна-Розена, определенное как пространственно-подобная связь между двумя вселенными с асимптотически минковскими областями пространства-времени в пределе больших расстояний от горизонтов.
	\item Рассмотрены соответствующие свойства симметрии и асимметрии обобщенного моста Эйнштейна-Розена на примерах метрик Рейсснера-Нордстрёма и Керра.
	\item С использованием диаграмм Картера-Пенроуза продемонстрированы свойства симметрии и асимметрии при движении пробного тела, например, космического корабля, сквозь множество различных вселенных внутри черной дыры.

%	\item На основе анализа \ldots
%	\item Численные исследования показали, что \ldots
%	\item Математическое моделирование показало \ldots
%	\item Для выполнения поставленных задач был создан \ldots
\end{enumerate}
